\documentclass[11pt]{article}
\usepackage[T1]{fontenc}
\usepackage[utf8]{inputenc}
\usepackage{lmodern}
\usepackage[margin=1in]{geometry}
\usepackage{amsmath,amssymb,amsthm,booktabs,microtype}
\usepackage{url}
\usepackage[hidelinks]{hyperref}
\hypersetup{pdftitle={Two-vertex cuts and a Lorentzian obstruction for bistable bipartite graphs},
  pdfauthor={Xiaojian Zeng}}
\newtheorem{theorem}{Theorem}
\newtheorem{lemma}[theorem]{Lemma}
\newtheorem{corollary}[theorem]{Corollary}
\DeclareMathOperator{\Hess}{Hess}
\DeclareMathOperator{\supp}{supp}
\newcommand{\N}{\mathbb{N}_0}
\newcommand{\R}{\mathbb{R}}
\title{Two-vertex cuts and a Lorentzian obstruction\newline
for bistable bipartite graphs}
\author{Xiaojian Zeng\thanks{PBC School of Finance, Tsinghua University,
100190. Email: \href{mailto:xiaojianzeng@foxmail.com}{xiaojianzeng@foxmail.com}.}}
\date{}

\begin{document}
\maketitle
\begin{abstract}
Let $F$ be a connected bipartite graph whose only maximum independent sets
are its two parts, each of size $k\ge2$. For the three-spin model in which
spins 1 and 2 must each occupy an independent set and spin 3 is unrestricted,
we express the coefficient of $x_1^{k-1}x_2^{k-1}x_3^2$ exactly in terms of
the connected components left by two-vertex deletions. The coefficient is
at least $2k^2$, with equality precisely when $F$ has no two-vertex cut.
The corresponding quadratic derivative has two positive Hessian eigenvalues
precisely when $F$ is not 3-connected. This gives one fixed antiferromagnetic
target witnessing non-Lorentzianity for all such sources of vertex
connectivity two, including every even cycle. For 3-connected sources the
same derivative has a zero eigenvalue and supplies no obstruction.
The general connected-clique characterization remains unresolved.
\end{abstract}

\section{Partition functions and the separator formula}

For a finite simple graph $F=(V,E)$ and a real symmetric matrix $A$ indexed
by a finite nonempty spin set $I$, define
\begin{equation}\label{eq:partition}
 Z_{F,A}(\mathbf{x})=
 \sum_{\sigma:V\to I}
 \left(\prod_{uv\in E}A_{\sigma(u),\sigma(v)}\right)
 \prod_{v\in V}x_{\sigma(v)}.
\end{equation}
Assignments are labeled, all vertices use the same variables, and there is
no factorial normalization. Following Lee, Oh and Seo~\cite{LOS}, call $A$
\emph{antiferromagnetic} if it is entrywise nonnegative and has at most one
positive eigenvalue, counted with multiplicity. Diagonal entries are allowed.
Their Theorem~3.1 proves that $Z_{K_n,A}$ is Lorentzian for every such $A$
and $n\ge2$; the case $n=1$ is immediate.
Their Conjecture~7.1 proposes a converse. We use the explicitly connected
formulation in Lee and Seo~\cite[Section~6]{LS}: a finite connected source
$F$ has Lorentzian $Z_{F,A}$ for every finite antiferromagnetic target $A$
if and only if $F$ is a clique. The result below treats a restricted class
in the unresolved direction of this characterization.

We recall the convention for Lorentzianity. A set $S\subseteq\N^I$ of
constant coordinate sum is \emph{M-convex} if, for $a,b\in S$ and
$a_i>b_i$, there is $j$ with $a_j<b_j$ and $a-e_i+e_j\in S$.
A nonzero homogeneous polynomial $f$ of degree $d\ge2$ with nonnegative
coefficients is Lorentzian precisely when $\supp(f)$ is M-convex and
\begin{equation}\label{eq:necessary}
 \Hess(\partial^\gamma f)\text{ has at most one positive eigenvalue}
 \quad\text{for every }\gamma\in\N^I,\quad |\gamma|=d-2.
\end{equation}
This is the characterization following Theorem~2.25 of
Br\"and\'en and Huh~\cite{BH}. The zero polynomial and nonnegative
homogeneous linear polynomials are included in the closed convention.
We will only need the necessary condition~\eqref{eq:necessary}.

A connected bipartite graph is \emph{bistable} when its two parts are its
only maximum independent sets, as in de Berg, Jansen and
Mukherjee~\cite{BJM}. The parts necessarily have the same size. Biconnectivity
of this class is known~\cite{BJM}; we rederive it below to make the proof
self-contained. A \emph{two-vertex cut} is an unordered pair whose deletion
disconnects the graph. Components consisting of one vertex are counted.

Throughout the theorem, fix
\begin{equation}\label{eq:target}
 A_\star=\begin{pmatrix}0&1&1\\1&0&1\\1&1&1\end{pmatrix}.
\end{equation}
Thus the first two spin classes are independent sets, while the third
class is unrestricted. Every allowed assignment has weight one.

\begin{theorem}\label{thm:main}
Let $F=(L\cup R,E)$ be a finite connected simple bipartite graph, with
$|L|=|R|=k\ge2$, whose only maximum independent sets are $L$ and $R$.
Put
\[
 f=Z_{F,A_\star},\qquad \beta=(k-1,k-1,0),\qquad
 c=[x_1^{k-1}x_2^{k-1}x_3^2]f.
\]
For distinct $u,v$, let $q(u,v)$ be the number of components of
$F-\{u,v\}$, and let
\[
 s=\#\bigl\{D\in\tbinom{L}{2}\cup\tbinom{R}{2}:F-D
       \text{ is disconnected}\bigr\}.
\]
Then
\begin{align}
 c&=\sum_{u\in L,\,v\in R}2^{q(u,v)}+2s,\label{eq:coefficient}\\
 \Hess(\partial^\beta f)
   &=2(k!)^2\begin{pmatrix}0&1&1\\1&0&1\\1&1&c/k^2\end{pmatrix}.
       \label{eq:hessian}
\end{align}
In particular, $F$ is 2-connected and $c\ge2k^2$, with equality exactly
when $F$ has no two-vertex cut. The matrix in~\eqref{eq:hessian} has
two positive eigenvalues exactly when $F$ is not 3-connected; in that
case $f$ is not Lorentzian. The target $A_\star$ is antiferromagnetic,
with spectrum $-1,1-\sqrt2,1+\sqrt2$.

If $F$ is 3-connected, this Hessian has one positive, one negative and
one zero eigenvalue. This last assertion does not imply that $f$ is Lorentzian.
\end{theorem}

The coefficient formula and its equality criterion connect a specific
quadratic derivative to vertex connectivity. They are the contribution
beyond the known clique sufficiency and bistability facts above. For
comparison, when $A$ is the adjacency matrix of a loopless complete graph,
\eqref{eq:partition} is a finite-variable specialization of Stanley's
chromatic symmetric function~\cite{Stanley}. Matherne, Morales and
Selover~\cite[Theorem~1.5]{MMS} prove Lorentzianity of these proper-coloring
polynomials for indifference graphs of abelian Dyck paths. Our target
allows adjacent vertices in spin 3, so their theorem concerns a different
polynomial. A related use of coloring coefficients to recover graph
structure appears in Zijian Zeng's reconstruction of proper trees with
multiplicity-isolated weighted cores~\cite[Theorem~1.1]{Zeng}: leading
star-basis data and the next coefficient layer recover the tree under
that hypothesis. Here a monomial coefficient of the three-spin partition
function detects two-vertex cuts under bistability. These comparisons
concern the cited sources and do not establish a global priority claim.

\section{Proof of the theorem}

Write $N(S)$ for the neighborhood of a vertex set in $F$.

\begin{lemma}\label{lem:expansion}
Under the hypotheses of Theorem~\ref{thm:main}, every nonempty proper
subset $S$ of either part satisfies
\begin{equation}\label{eq:expansion}
 |N(S)|\ge |S|+1.
\end{equation}
Consequently, every vertex has degree at least two, and $F$ is 2-connected.
\end{lemma}
\begin{proof}
Suppose $\varnothing\ne S\subsetneq L$ and $|N(S)|\le |S|$.
Then $S\cup(R\setminus N(S))$ is independent and has size at least $k$.
A larger size contradicts the maximum size of the two parts. Equality gives
a maximum independent set distinct from both $L$ and $R$, again a
contradiction. Interchanging the parts proves~\eqref{eq:expansion} on
both sides. Applying it to singletons gives the degree bound.

Delete $u\in L$. Each component of $F-u$ has a nonempty $R$-part:
any remaining $L$-vertex retains all its neighbors. If there were several
components, write their parts as $L_i,R_i$. Each $R_i$ is then a nonempty
proper subset of $R$, and
\[
 |R_i|+1\le |N(R_i)|\le |L_i|+1.
\]
Summing would give $k\le k-1$. Thus $F-u$ is connected. The argument
for $u\in R$ is symmetric. Since $|V|\ge4$, the graph is 2-connected.
\end{proof}

\begin{proof}[Proof of Theorem~\ref{thm:main}]
To count $c$, choose the unordered pair $D$ assigned spin 3. The other
vertices must be properly colored by spins 1 and 2, with $k-1$ vertices
of each color. We distinguish the two possible positions of $D$.

If $D=\{u,v\}$ with $u\in L$ and $v\in R$, every remaining vertex
retains a neighbor by Lemma~\ref{lem:expansion}. Each component of $F-D$
therefore has nonempty parts $L_i,R_i$, both proper in the original parts.
Its neighborhoods in $F$ lie respectively in $R_i\cup\{v\}$ and
$L_i\cup\{u\}$. Applying~\eqref{eq:expansion} gives
\[
 |L_i|+1\le |R_i|+1,\qquad |R_i|+1\le |L_i|+1.
\]
Thus every component is balanced. Its two color orientations are
independent of the other components, and either orientation preserves
the required totals. This pair contributes $2^{q(u,v)}$.

Next suppose $D\subset L$. If $F-D$ is connected, its bipartition
has sizes $k-2,k$ and it has no balanced proper coloring. If it is
disconnected, each component has a nonempty proper $R$-part $R_i$.
Since $N(R_i)\subseteq L_i\cup D$, expansion gives
$|R_i|\le |L_i|+1$. When $L_i\ne\varnothing$, it is a proper
subset of $L$ and $N(L_i)\subseteq R_i$, so expansion also gives
$|R_i|\ge |L_i|+1$. When $L_i=\varnothing$, the component is
one isolated $R$-vertex, and the same equality holds. Hence every
component satisfies
\[
 |R_i|-|L_i|=1.
\]
The total imbalance is two, so there are exactly two components.
Exactly two of their four orientation choices cancel the imbalances
and yield $k-1$ vertices of each color. Thus each such cut contributes
two. Deletion within $R$ is symmetric. This proves~\eqref{eq:coefficient},
including singleton components and the case $k=2$.

Every remainder has at least two vertices, so $q(u,v)\ge1$.
Each of the $k^2$ cross-part terms is at least two, and a term exceeds
two precisely for a cross-part cut. The term $2s$ is positive precisely
when there is a same-part cut. This proves $c\ge2k^2$ and its equality
criterion. By Lemma~\ref{lem:expansion}, absence of a two-vertex cut is
equivalent here to 3-connectivity.

It remains to compute the Hessian without changing normalization.
The independence number is $k$, so the coefficients with exponents
$(k+1,k-1,0)$ and $(k-1,k+1,0)$ vanish. Connectedness gives
\[
 [x_1^kx_2^k]f=2.
\]
If the spin-1 class has size $k$, it must be $L$ or $R$. Choosing
one of the $k$ vertices of the other part for spin 3 gives
\[
 [x_1^kx_2^{k-1}x_3]f=2k,
 \qquad [x_1^{k-1}x_2^kx_3]f=2k.
\]
For $\alpha\in\N^3$, write $\alpha!=\alpha_1!\alpha_2!\alpha_3!$.
Since $|\beta|=2k-2$, the constant $(i,j)$ entry of
$\Hess(\partial^\beta f)$ is
\[
 (\beta+e_i+e_j)!\,[\mathbf{x}^{\beta+e_i+e_j}]f.
\]
These coefficients give~\eqref{eq:hessian}; in particular, the $(3,3)$
entry is $2((k-1)!)^2c$, with the factor two coming from $2!$.

Set $t=c/k^2\ge2$ and divide the Hessian by the positive number
$2(k!)^2$. The vector $(1,-1,0)$ is an eigenvector of eigenvalue $-1$.
On its orthogonal complement, in the orthonormal basis
$(1,1,0)/\sqrt2,(0,0,1)$, the matrix is
\[
 B_t=\begin{pmatrix}1&\sqrt2\\\sqrt2&t\end{pmatrix},
 \qquad \det B_t=t-2,\quad \operatorname{tr}B_t=1+t.
\]
For $t>2$, both eigenvalues of $B_t$ are positive. For $t=2$, they
are $3$ and $0$. This proves the asserted Hessian inertias. Finally,
\[
 \det(\lambda I-A_\star)=(\lambda+1)(\lambda^2-2\lambda-1),
\]
so the target is entrywise nonnegative with exactly one positive
eigenvalue. If $F$ has a two-vertex cut, the order-$2k-2$ derivative
just computed violates~\eqref{eq:necessary}, proving non-Lorentzianity.
\end{proof}

\section{Even cycles and the equality boundary}

\begin{corollary}\label{cor:examples}
For $F=C_{2k}$, $k\ge2$, the coefficient in Theorem~\ref{thm:main} is
$c=6k(k-1)$, and $Z_{F,A_\star}$ is not Lorentzian.
For $F=K_{k,k}$, $k\ge3$, it is $c=2k^2$.
For $K_{2,2}=C_4$ it is instead $c=12$.
\end{corollary}
\begin{proof}
An independent set in $C_{2k}$ has at most $k$ vertices. If it has
$k$, its $k$ cyclic gaps all have length at least two and sum to $2k$,
so each gap has length two. The maximum independent sets are therefore
the two parts. Of the $k^2$ cross-part pairs, $2k$ are adjacent and
leave one component; each of the other $k^2-2k$ pairs leaves two.
All $2\binom{k}{2}=k(k-1)$ same-part pairs are cuts. Thus
\[
 c=4k+4(k^2-2k)+2k(k-1)=6k(k-1)>2k^2.
\]
In $K_{k,k}$ every independent set lies in one part. For $k\ge3$,
deleting any two vertices leaves a connected complete bipartite graph,
so Theorem~\ref{thm:main} gives equality. For $k=2$, the two same-part
pairs each leave two isolated vertices, giving $c=8+4=12$.
\end{proof}

Equality is a limitation of the selected derivative, not a criterion for
Lorentzianity of the partition function. Other derivatives, or other
antiferromagnetic targets, may still provide an obstruction on a
3-connected bistable source. The theorem makes no assertion for arbitrary
bipartite sources lacking the maximum-independent-set hypothesis, and it
does not resolve the connected-clique characterization.

\section{Exact computational checks}

The proof above is independent of computation. As a finite check, two
separately implemented Python programs examined every labeled bipartite
graph with fixed parts of size $k=2,3,4$. All $2^{k^2}$ edge sets were
considered, with minimum degree two used only as the necessary prefilter
proved in Lemma~\ref{lem:expansion}. The first program tests all vertex
subsets for the independent-set hypothesis and counts disjoint independent
color classes. The second tests all $k$-subsets and counts assignments
directly from $A_\star$ by vertex recursion; its component computation
uses a separate breadth-first search. Testing $k$-subsets suffices because
an independent set larger than $k$ contains a $k$-subset different from
both parts. The computations use exact integer or rational arithmetic.

The programs agree on all six coefficients needed for the Hessian,
the separator identity and its strict/equality criterion. The exhaustive
coverage is as follows; graphs are counted with their fixed labels, without
quotienting by isomorphism.
\begin{center}
\begin{tabular}{rrrrr}
\toprule
$k$ & Edge sets & Bistable graphs & $c>2k^2$ & $c=2k^2$\\
\midrule
2 & 16 & 1 & 1 & 0\\
3 & 512 & 34 & 33 & 1\\
4 & 65\,536 & 6\,785 & 6\,576 & 209\\
\bottomrule
\end{tabular}
\end{center}
There are 6\,820 qualifying graphs in total. A further 22 structured
records through $k=7$ and 15 representative records from the exhaustive
domain were checked from explicit edge lists by the second program.
These records are not all distinct and do not constitute exhaustive
larger-order coverage. The cube, represented as $K_{4,4}$ minus a perfect
matching, is a checked equality example with $c=32$.
The accompanying computational dossier contains the programs and exact
outputs. The computations corroborate the all-$k$ proof; the eigenvalue
signs used in the theorem follow from the displayed $2\times2$ reduction.

\paragraph{Computational and AI assistance.}
OpenAI's Codex was used to develop the counting argument, generate Python
verification programs, inspect primary literature and prepare the manuscript.
A separate model context reconstructed the argument and implemented the
second finite verifier. This was an additional checking pass, not a
formal proof-assistant verification or journal peer review. The mathematical
justification is the explicit proof above, with computation serving as
finite corroboration.

\bibliographystyle{amsplain}
\bibliography{references}
\end{document}
